\documentclass[10pt,letterpaper]{article}
\usepackage{cornell-notes}

\setDocTitle{CISSP CBK Cornell Notes}
\setDocSubtitle{Domain 3: Security Architecture and Engineering}
\setDocVersion{Sixth Edition \textbullet\ 2021}
\setDocStatus{Study Companion}
\setDocOwner{Security Certification Study Collection}
\setDocAuthor{Arthur Deane and Aaron Kraus}
\setDocDate{\today}
\setDocSummary{Cornell-format study and review notes for Domain 3 of the CISSP CBK reference.}

\setCornellCollection{CISSP CBK Cornell Notes}
\setCornellUnitType{Domain}
\setCornellUnitNumber{3}
\setCornellUnitTitle{Security Architecture and Engineering}
\setCornellSourceLabel{CBK}
\setCornellSourceTitle{THE OFFICIAL (ISC)\textsuperscript{2} CISSP CBK REFERENCE}
\setCornellSourceAuthor{Arthur Deane and Aaron Kraus}
\setCornellSourceEdition{Sixth Edition \textbullet\ 2021}
\setCornellSourceRange{pp. 147--281}
\setCornellCompanionLabel{Study and review companion}

\begin{document}
\makecornelltitle

\section*{Domain Overview}
Domain 3 turns security requirements into architecture. It links secure design principles, formal security models, system/hardware capabilities, modern compute architectures, cryptographic engineering, cryptanalytic threats, and physical facility design. The central idea is assurance: design multiple mediated boundaries, minimize privilege and attack surface, choose controls from requirements, and validate that the complete system preserves its intended security properties.

\begin{summarybox}{Review Objectives}
Explain ISO/IEC 19249 and related secure-design principles; compare Bell-LaPadula, Biba, Clark-Wilson, Brewer-Nash, Take-Grant, information-flow, and noninterference models; select system controls; evaluate modern architectures; choose cryptographic mechanisms and lifecycle practices; recognize cryptanalytic attacks; and apply site, data-center, environmental, and fire-safety controls.
\end{summarybox}

\section{Secure Engineering and Design Principles}
\source{pp. 149--167}
\begin{cornell}
\noterow{ISO/IEC 19249 purpose}{Catalog architectural and design principles for secure products, systems, and applications. The chapter presents five architectural principles and five design principles that help structure security before implementation details are chosen.}
\noterow{Five architectural principles}{\textbf{Domain separation, layering, encapsulation, redundancy, virtualization}. Together they create boundaries, abstractions, fault tolerance, and isolation that make security policies easier to enforce and analyze.}
\noterow{Domain separation}{Group components with similar security attributes into domains and force cross-domain communication through well-defined, fully mediated channels. Network segmentation and classification-based data domains are practical examples.}
\noterow{Layering}{Build hierarchical abstractions in which upper layers use controlled services of lower layers. Security boundaries at each layer only work if higher layers cannot bypass the intermediate enforcement point. Layering also simplifies validation.}
\noterow{Encapsulation}{Hide an object's internal implementation and require access through controlled functions/interfaces. If all operations are mediated by the encapsulation interface, policy can be checked consistently and internal state is less exposed.}
\noterow{Redundancy}{Replicate components so failure or overload does not destroy availability, and maintain consistent state across replicas. Detect failure, fail over safely, and where a transaction is partially complete, roll back or reconcile state rather than duplicate/corrupt it.}
\noterow{Virtualization}{Abstract underlying compute/storage/network resources so multiple isolated environments can share hardware. Security depends on strong isolation, hypervisor/container controls, resource management, and protection of the management plane.}
\noterow{Five ISO design principles}{\textbf{Least privilege, attack-surface minimization, centralized parameter validation, centralized general security services, preparing for error and exception handling}. These reduce unnecessary authority and duplicated security logic while making failure behavior explicit.}
\noterow{Least privilege}{Give a subject/process only the permissions necessary, only for the needed purpose and duration. Just-in-time elevation complements least privilege by shrinking the time privileged access exists.}
\noterow{Attack-surface minimization}{Remove/disable unnecessary services, ports, features, accounts, interfaces, and complexity. Every exposed mechanism is an opportunity for defect or misuse; hardening reduces the number of opportunities an attacker can reach.}
\noterow{Centralized parameter validation}{Use well-reviewed validation mechanisms for untrusted or cross-boundary input instead of many inconsistent implementations. Validate type, length, range, format, and allowed values before data reaches sensitive operations.}
\noterow{Centralized security services}{Centralize common functions such as authentication, authorization, logging, and key management when doing so improves consistency and reviewability. Avoid creating an unprotected single point of failure; design the service itself for resilience and trust.}
\noterow{Errors and exceptions}{Fail in a predictable, secure state; do not leak sensitive implementation details in error messages; log enough for diagnosis; and define recovery/rollback behavior. Security must remain intact when normal processing breaks.}
\noterow{Threat modeling}{Apply structured threat identification during design, before implementation makes changes expensive. Map trust boundaries, assets, entry points, adversary actions, and mitigations; revisit the model when architecture or threats change.}
\noterow{Secure defaults}{Default configuration should be the safer state; require deliberate action to weaken protection or enable risky functionality. This protects users who never change defaults and reduces accidental exposure.}
\noterow{Fail securely}{Failure should preserve required confidentiality, integrity, and authorization. A fail-open design can favor availability but expose protected resources; fail-closed can protect access but reduce availability. Select behavior from safety/business requirements.}
\noterow{Separation of duties}{Split sensitive functions so one person/process cannot complete a high-risk action alone. This limits fraud, mistakes, and abuse; dual control and two-person integrity are common implementations.}
\noterow{Keep it simple}{Complex designs are harder to understand, verify, patch, and operate. Minimize unnecessary mechanisms and dependencies so security behavior is reviewable and predictable.}
\noterow{Trust, but verify}{Do not rely on asserted trust without evidence. Authenticate, authorize, log, monitor, attest, test, and periodically reassess trusted relationships and components.}
\noterow{Zero trust}{Treat network location as insufficient evidence of trust. Continuously evaluate identity, device/context, least privilege, and policy for access to resources; authenticate/authorize communications and protect them regardless of request origin.}
\noterow{Privacy by Design}{Embed privacy across the lifecycle rather than bolt it on. The seven principles emphasize proactive prevention, privacy by default, embedded privacy, positive-sum functionality, lifecycle protection, visibility/transparency, and user-centric respect.}
\noterow{Shared responsibility}{In cloud services the provider and customer divide security responsibilities. The split changes across IaaS/PaaS/SaaS and provider contracts; customers remain responsible for their assigned configuration, identities, data, and use rather than assuming the provider owns all security.}
\noterow{Defense in depth}{Combine distinct layers/controls so a single failure does not immediately cause compromise. Assume breach, include prevention plus detection/response, and avoid hidden common-mode failures that make nominally redundant controls fail together.}
\end{cornell}

\section{Fundamental Security Models}
\source{pp. 168--175}
\begin{cornell}
\noterow{Why security models?}{A model abstracts a system and expresses security policy as rules that can be reasoned about or formally verified. Models reduce ambiguity about what the architecture must preserve.}
\noterow{State machines and lattices}{A finite-state model analyzes whether transitions keep the system in secure states. A lattice assigns ordered security levels to subjects/objects and rules governing allowed operations based on those levels.}
\noterow{Information-flow model}{Controls where information may move among objects/systems based on classification or other policy. Bell-LaPadula and Biba are classic information-flow models with different primary objectives.}
\noterow{Noninterference}{Actions at a higher security level should not cause observable effects at a lower level that reveal higher-level information. The goal is to prevent unintended information leakage through interaction between security domains.}
\noterow{Bell-LaPadula (BLP)}{\textbf{Confidentiality}. Simple Security Property: \textbf{no read up}. Star Property: \textbf{no write down}. A discretionary property uses an access matrix. The model prevents lower-cleared subjects from reading higher data and prevents leakage by writing sensitive data to lower levels.}
\noterow{Biba}{\textbf{Integrity}. Conceptually reverses BLP information flow: \textbf{no read down} protects a high-integrity subject from lower-integrity input; \textbf{no write up} prevents lower-integrity subjects from contaminating higher-integrity objects.}
\noterow{Clark-Wilson}{Commercial integrity model built around well-formed transactions, separation of duties, and controlled transformation procedures. Users act through authorized programs rather than directly manipulating constrained data items; integrity verification checks valid state.}
\noterow{Brewer-Nash / Chinese Wall}{Dynamic confidentiality model that prevents conflicts of interest. Access to one company's sensitive data can restrict subsequent access to a competitor's sensitive data within the same conflict class. Permissions depend on prior accesses, not only static clearance.}
\noterow{Take-Grant}{Graph-based model representing subjects/objects and rights. Take and grant rules describe how permissions can propagate. It is useful for reasoning about whether a subject can eventually acquire a right through allowed transformations.}
\end{cornell}

\begin{exambox}{Classic Model Memory Cue}
Bell-LaPadula protects \textbf{confidentiality}: no read up, no write down. Biba protects \textbf{integrity}: no read down, no write up. Clark-Wilson protects commercial integrity through authorized transformations and separation of duties.
\end{exambox}

\section{Controls and Information-System Security Capabilities}
\source{pp. 175--186}
\begin{cornell}
\noterow{Select controls from requirements}{Start with business/security requirements, risk, classification, law/contract, and architecture. Select a baseline, tailor it, document rationale, and revisit controls after incidents, major system/organizational changes, new threats, or framework updates.}
\noterow{Memory protection}{Operating systems and processors isolate process memory so one program cannot directly access another's assigned regions. Dual-mode operation separates privileged/kernel instructions from user mode. Protection depends on both correct hardware and operating-system behavior.}
\noterow{ASLR}{Address Space Layout Randomization makes important code/data addresses less predictable, increasing difficulty for exploits that depend on known memory locations. It is defense in depth, not a substitute for memory-safe design and patching.}
\noterow{Secure cryptoprocessor}{Tamper-resistant hardware with a constrained interface for key generation/storage, random numbers, encryption/signing, trusted boot, and acceleration. Limiting exposed functions makes assurance easier and keeps sensitive key material away from general-purpose software.}
\noterow{TPM}{Trusted Platform Module is a standardized secure cryptoprocessor. Core concepts include \textbf{attestation} of known state, \textbf{binding} data/keys to a platform, and \textbf{sealing} data so it can be released only when platform measurements match required state.}
\noterow{HSM}{Hardware Security Module is dedicated tamper-resistant hardware for enterprise cryptographic key generation, storage, and operations. It centralizes protected key use and often supports high throughput, auditing, role separation, and strong administrative controls.}
\end{cornell}

\section{Architecture and Solution-Element Vulnerabilities}
\source{pp. 187--223}
\begin{cornell}
\noterow{Client-based systems}{Endpoints face malicious content, local privilege abuse, stolen credentials, vulnerable applications, and physical loss. Harden, patch, encrypt, restrict privilege, protect credentials, use endpoint controls, and design assuming a client can be hostile.}
\noterow{Server-based systems}{Servers concentrate services/data and are attractive targets. Minimize exposed services, segment management interfaces, patch/harden, protect privileged access, log/monitor, isolate workloads, and design redundancy/recovery according to business impact.}
\noterow{Database systems}{Enforce database authorization and least privilege; separate application/service/admin roles; validate queries/input; encrypt sensitive data and backups; audit privileged/data access; patch the DBMS; and protect integrity, concurrency, and recovery mechanisms.}
\noterow{Cryptographic-system weaknesses}{Risk comes from weak/deprecated algorithms, bad randomness, key exposure/reuse, unsafe modes/protocols, poor implementation, insecure defaults, side channels, and failed certificate/key lifecycle. Correct primitives do not rescue an incorrect cryptosystem.}
\noterow{Industrial control systems (ICS)}{ICS/SCADA/OT controls physical processes, so safety and availability can dominate. Legacy protocols, long lifecycles, fragile equipment, flat networks, vendor remote access, and convergence with IP networks create risk. Segment, tightly mediate access, monitor, and test changes carefully.}
\noterow{Cloud service models}{\textbf{IaaS}: customer manages more of OS, applications, identities/data. \textbf{PaaS}: provider additionally manages platform/runtime. \textbf{SaaS}: provider manages the application stack while customer still manages users, data, configuration, and use. Verify the exact contract/model.}
\noterow{Cloud deployment/security}{Public, private, hybrid, and community approaches change ownership and exposure, not the need for risk management. Evaluate tenancy/isolation, data location, identity, APIs, logging, encryption/key control, provider assurance, resilience, and exit/portability.}
\noterow{Distributed systems}{Components communicate across failure and trust boundaries. Secure service identity, network communication, consistency, authorization, time/synchronization assumptions, distributed logging, and failure handling; account for a node or link becoming unavailable or malicious.}
\noterow{Internet of Things (IoT)}{Resource constraints, default credentials, weak update mechanisms, insecure services, long lifecycles, physical exposure, and large device populations enlarge attack surface. Inventory, change defaults, segment, authenticate, encrypt, patch, disable unused services, and design secure update/recovery.}
\noterow{Microservices}{Small independently deployable services communicate through APIs. Benefits include isolation and independent scaling/deployment, but many service identities, APIs, secrets, dependencies, and network paths increase management complexity. Authenticate/authorize service calls and protect API boundaries.}
\noterow{Containers}{Package applications with dependencies while sharing the host kernel. Protect image provenance, registry and orchestration control planes, secrets, runtime privileges, namespaces/cgroups, host kernel, networking, and patch lifecycle. Container isolation differs from a separate-OS VM boundary.}
\noterow{Serverless}{Provider dynamically runs short-lived functions/resources. Granular functions can reduce persistent server management, but risks include broken authentication, over-privileged roles, insecure dependencies, event/data injection, secrets, monitoring gaps, and provider/platform dependency.}
\noterow{Embedded systems}{Dedicated computing built into devices often has constrained resources, long replacement cycles, and physical exposure. Use secure boot/firmware, signed updates, hardware roots of trust, minimal interfaces, unique credentials, and tamper resistance where risk justifies it.}
\noterow{High-performance computing (HPC)}{Clusters/supercomputers aggregate valuable compute and data and often emphasize performance. Secure management/scheduling, high-speed interconnects, shared storage, privileged tooling, identities, research data, and availability without assuming specialized systems are immune to ordinary OS/software vulnerabilities.}
\noterow{Edge computing}{Moves compute/storage near data sources/users, reducing latency and central bandwidth. Distributed, remotely managed, and physically exposed nodes create inventory, patching, tamper, identity, connectivity, and failover challenges. Protect local data and secure central management.}
\noterow{Virtualized systems}{Hypervisors mediate shared hardware among VMs. Type 1 runs directly on hardware; Type 2 runs atop a host OS. Secure the hypervisor/management plane, templates/snapshots, virtual networks, resource allocation, VM migration, and isolation; prevent VM escape and cross-tenant leakage.}
\end{cornell}

\section{Cryptographic Solutions}
\source{pp. 224--256}
\begin{cornell}
\noterow{What security properties can crypto provide?}{Depending on mechanism and key management: confidentiality, integrity/tamper detection, authentication, and nonrepudiation. Cryptography cannot by itself correct bad authorization, compromised endpoints, weak operations, or lost keys.}
\noterow{Core terms}{\textbf{Plaintext} is readable input; \textbf{ciphertext} is encrypted output; a \textbf{cipher/algorithm} defines transformation; a \textbf{key} selects the specific transformation; \textbf{encryption/decryption} transform between plaintext and ciphertext.}
\noterow{Cryptographic lifecycle}{Select appropriate function/algorithm/key length/mode; generate keys with strong entropy; distribute/store/use them securely; rotate/renew as required; revoke/expire compromised or obsolete material; archive/recover where authorized; and destroy keys when retention ends. Plan algorithm agility.}
\noterow{Algorithm/key selection}{Use well-vetted current standards rather than custom cryptography. Consider threat model, sensitivity, required lifetime, performance, hardware support, regulatory guidance, mode/protocol, and expected advances in cryptanalysis/computing. Longer keys increase brute-force resistance but can cost performance.}
\noterow{Symmetric cryptography}{One shared secret key performs encryption/decryption. It is fast and well suited to bulk data at rest or in transit, but secure key distribution and protection become difficult as the number of parties grows. Block and stream ciphers process data differently.}
\noterow{AES}{Advanced Encryption Standard (Rijndael) is the chapter's current recommended symmetric standard. It uses a 128-bit block and 128-, 192-, or 256-bit keys. AES replaced DES; key length and safe mode/protocol selection both matter.}
\noterow{DES / 3DES / RC cue}{DES's 56-bit key is obsolete. 3DES extended DES by multiple passes but is legacy/deprecated for new designs. Rivest ciphers include several historical algorithms; RC4 in particular is unsafe for modern use. Recognize legacy use without selecting it for new protection.}
\noterow{Block modes}{ECB encrypts blocks independently and can reveal patterns; chained/feedback modes use prior state or feedback to hide repeated structure. Modes require correct initialization vectors/nonces and integrity protection where appropriate. Prefer standardized authenticated constructions/protocols rather than ad hoc composition.}
\noterow{Asymmetric/public-key cryptography}{Uses a mathematically related public/private key pair. Public keys can be shared; private keys must remain secret. It solves important key-distribution/signature problems but is computationally heavier, so hybrid protocols commonly use asymmetric crypto to establish/authenticate symmetric session keys.}
\noterow{Diffie-Hellman-Merkle}{Key-agreement method enabling parties to establish a shared secret over an insecure channel. It is not itself bulk encryption. Without authentication, key exchange can be vulnerable to a man-in-the-middle, so protocols bind the exchange to authenticated identities.}
\noterow{RSA}{Public-key algorithm based on difficulty of factoring large integers; used for encryption/signature functions in appropriate schemes. Security depends on key size, padding/protocol, random generation, private-key protection, and implementation.}
\noterow{ElGamal and ECC}{ElGamal is based on the discrete logarithm problem and supports key exchange/signature-related uses. Elliptic Curve Cryptography provides comparable security with substantially smaller keys; related schemes include ECDH for key agreement and ECDSA for signatures.}
\noterow{Hashing}{One-way function maps arbitrary input to fixed-length output. Secure hashes support integrity checking, signatures/MAC constructions, and password-storage schemes. Collision/preimage resistance matters; obsolete hashes should not be selected for security-sensitive new designs.}
\noterow{Password hashing}{Use a deliberately expensive password-based KDF with a unique random salt per credential/change. Salting defeats one precomputed table from working across all users; a separately protected pepper can add another secret, but it does not replace slow hashing or MFA.}
\noterow{PKI}{Public Key Infrastructure binds public keys to identities using certificates and trust relationships. Certification Authorities (CAs) issue/sign certificates; relying parties validate chain, identity/name, time validity, usage constraints, and revocation/status before trusting a key.}
\noterow{Certificate lifecycle}{Enrollment/identity proofing, issuance, distribution/use, renewal, revocation, and expiration must be managed. Protect CA and subject private keys; a compromised CA can undermine trust far beyond one certificate. CRLs/online status mechanisms help identify revoked certificates.}
\noterow{Key management}{Generate with strong randomness; distribute securely; restrict access; store private/symmetric keys in protected hardware where appropriate; rotate/expire; escrow/recover only under authorized policy; revoke compromised keys; and destroy keys when data no longer requires access.}
\noterow{Digital signature}{Hash the message and use the signer's private-key operation to produce a signature; verification uses the corresponding public key and comparison to the message hash. This supports integrity, origin authentication, and nonrepudiation when identity/key trust is sound. It does not provide confidentiality by itself.}
\noterow{Digital certificate}{Signed data structure binding a subject identity/name to a public key plus metadata such as issuer, validity period, serial number, algorithm, and usage. Trust depends on validating the chain and certificate conditions, not merely seeing that a certificate exists.}
\noterow{Nonrepudiation and integrity}{Nonrepudiation requires evidence that makes denial of an action difficult; digital signatures and trustworthy identity/key handling contribute. Integrity can also be detected with hashes/MACs, but shared-secret MACs do not provide the same third-party nonrepudiation as a properly managed signature.}
\end{cornell}

\begin{exambox}{Crypto Selection Cue}
Bulk confidentiality normally favors symmetric crypto. Public-key techniques establish identity/trust and help with key establishment/signatures. Hashes detect change and protect password verifiers when used in an appropriate slow, salted construction. A digital signature provides integrity/authentication/nonrepudiation -- not secrecy.
\end{exambox}

\section{Cryptanalytic Attacks}
\source{pp. 257--264}
\begin{cornell}
\noterow{Brute force}{Try possible keys/credentials until one works. Defenses include adequate key length and slow salted password KDFs; rate limiting/MFA protect online authentication. Computational feasibility changes over time.}
\noterow{Ciphertext-only}{Attacker observes only ciphertext and uses statistical/structural properties to infer plaintext or keys. Strong modern algorithms/modes aim to remain secure in this model.}
\noterow{Known plaintext}{Attacker knows corresponding plaintext and ciphertext for some messages. Secure algorithms must withstand this realistic condition because protocol headers and predictable content are common.}
\noterow{Chosen plaintext}{Attacker can choose plaintext and observe corresponding ciphertext. Modern constructions are designed against such oracle-like capabilities when used correctly.}
\noterow{Frequency analysis}{Uses symbol frequency to attack simple substitution and other weak classical ciphers. Strong modern encryption diffuses statistical patterns so direct frequency analysis is ineffective.}
\noterow{Chosen ciphertext}{Attacker can submit chosen ciphertext and learn the corresponding plaintext/error behavior. Padding/oracle flaws show why decryption interfaces must not reveal exploitable distinctions.}
\noterow{Implementation attacks}{Exploit defects in hardware/software/protocol implementation rather than breaking the mathematical primitive. Correct algorithms can fail through bad randomness, memory bugs, validation errors, unsafe APIs, or configuration.}
\noterow{Side-channel attacks}{Infer secrets from timing, cache behavior, power consumption, electromagnetic emanations, or error information. Defenses can include constant-time operations, isolation, shielding, blinding, noise, and tamper-resistant hardware.}
\noterow{Fault injection}{Induce voltage, clock, radiation, temperature, or other faults and analyze incorrect outputs/states to recover secrets or bypass checks. Tamper-resistant design and fault detection/redundant verification mitigate.}
\noterow{Timing attacks}{Exploit data/key-dependent execution duration. Implement sensitive operations so runtime and observable behavior are as independent of secrets as practical.}
\noterow{Man-in-the-middle}{Attacker interposes between parties and may relay/alter key establishment. Authenticate peers and bind key exchange to validated identities/certificates.}
\noterow{Pass the hash / Kerberos exploitation}{Authentication material itself can become a reusable credential. Protect credential stores and tickets/hashes, restrict privilege, isolate admin sessions, monitor anomalous authentication, and use modern credential protections.}
\noterow{Ransomware}{Malware encrypts or denies access to victim data for extortion. Crypto is used by the attacker rather than broken. Defense includes prevention, least privilege, segmentation, detection, tested offline/immutable backups, incident response, and recovery readiness.}
\end{cornell}

\section{Site and Facility Security}
\source{pp. 265--280}
\begin{cornell}
\noterow{Site selection}{Assess natural hazards, crime/civil unrest, neighboring facilities, transportation/flight paths, utilities, communications, emergency response, environmental conditions, jurisdiction, and geographic correlation with alternate sites. Physical design should follow risk, not convenience alone.}
\noterow{Defense in depth for facilities}{Use property/perimeter barriers, lighting/surveillance, controlled entrances, guards/reception, badges/biometrics where appropriate, visitor controls, internal zones, locked rooms/racks, and logging. Place increasingly sensitive assets behind progressively stronger boundaries.}
\noterow{Wiring closets / IDFs}{Protect intermediate distribution facilities because unauthorized physical access can expose network cabling and equipment. Lock, monitor, control keys/badges, avoid using them as general storage, and restrict access in multi-tenant/colocation facilities.}
\noterow{Server rooms / data centers}{Control entry, visitors, racks, cabling, media, environmental systems, and management interfaces. Design redundancy for power, cooling, network, and processing based on availability requirements; test the redundancy rather than assuming it works.}
\noterow{Media storage}{Use access controls plus stable temperature/humidity, filtration, fire/water protection, appropriate shelving/containers, and location separation. Consider media degradation/obsolescence and acclimatization when moving media between environments.}
\noterow{Evidence storage}{Preserve chain of custody with indelible in/out logs, strict access policy, surveillance, and dual controls/locked inner storage where needed. Every handling event must be attributable and evidence protected from tampering/contamination.}
\noterow{Restricted and work areas}{Separate public, controlled, and restricted zones; protect sensitive screens/documents, prevent tailgating, use visitor escort and clean-desk practices as required, and place assets so casual/unauthorized observation or access is difficult.}
\noterow{Power resilience}{Use appropriate feeds, UPS systems, surge/power conditioning, and generators. UPS bridges short interruptions and generator startup; capacity, fuel, maintenance, load testing, alerting, and failure modes all matter.}
\noterow{HVAC}{Availability depends on temperature, airflow, humidity, and filtration. Low humidity increases static risk; high humidity can cause condensation. Provide monitoring, alerting, maintenance, and redundant cooling consistent with the facility's required uptime.}
\noterow{Environmental threats}{Consider fire, flood, storms, earthquake, external hazards, burst pipes, roof/water leaks, dust/contaminants, and correlated events. Keep water sources away from critical equipment when practical and ensure alternate sites are not exposed to the same regional hazard.}
\noterow{Fire tetrahedron}{Fire requires heat, fuel, oxygen, and a sustaining chemical reaction. Prevention/suppression removes or interrupts one or more components; choose the agent for the fuel and occupied-space constraints.}
\noterow{Sprinkler/gaseous systems}{Wet-pipe sprinklers are common/reliable but water can damage equipment. Dry-pipe/deluge/pre-action designs fit specialized hazards. Gaseous agents can avoid water damage but require enclosure/HVAC coordination and life-safety precautions. Halon is legacy and environmentally harmful.}
\noterow{Fire classes}{\textbf{A}: ordinary solid combustibles; \textbf{B}: flammable liquids/gases; \textbf{C}: energized electrical equipment; \textbf{D}: combustible metals; \textbf{F/K}: cooking oils/greases. Extinguishers must match the fire class, and personnel need training/evacuation priorities.}
\end{cornell}

\section{Critical Comparisons}
\begin{center}\small
\begin{tabularx}{\textwidth}{>{\bfseries\color{Navy}\raggedright\arraybackslash}p{0.22\textwidth}>{\raggedright\arraybackslash}X>{\raggedright\arraybackslash}X}
\toprule Concept & First idea & Distinguishing idea\\\midrule
Bell-LaPadula / Biba & Confidentiality: no read up, no write down & Integrity: no read down, no write up\\
Symmetric / asymmetric & One shared secret; fast bulk encryption & Public/private pair; key establishment, encryption/signatures\\
Hash / encryption & One-way digest for integrity/verifiers & Reversible with key for confidentiality\\
Signature / encryption & Integrity, origin authentication, nonrepudiation & Confidentiality when properly applied\\
VM / container & Separate guest OS through hypervisor & Shared host kernel with process isolation\\
IaaS / PaaS / SaaS & Customer manages most of stack & Provider responsibility increases toward SaaS\\
Type 1 / Type 2 hypervisor & Runs on hardware & Runs on a host operating system\\
\bottomrule
\end{tabularx}
\end{center}

\section{Key Terms and Acronyms}
\begin{summarybox}{Rapid-Review Glossary}
\textbf{AES} -- Advanced Encryption Standard \quad
\textbf{ASLR} -- Address Space Layout Randomization \quad
\textbf{BLP} -- Bell-LaPadula \quad
\textbf{CA} -- Certification Authority \quad
\textbf{ECC} -- Elliptic Curve Cryptography \quad
\textbf{HSM} -- Hardware Security Module \quad
\textbf{ICS} -- Industrial Control System \quad
\textbf{IoT} -- Internet of Things \quad
\textbf{PKI} -- Public Key Infrastructure \quad
\textbf{TPM} -- Trusted Platform Module \quad
\textbf{VM} -- Virtual Machine \quad
\textbf{Zero trust} -- resource access without implicit trust based on network location.
\end{summarybox}

\section{High-Yield Review}
\begin{itemize}
\item ISO 19249 architecture: domain separation, layering, encapsulation, redundancy, virtualization.
\item ISO 19249 design: least privilege, minimize attack surface, centralize validation/security services, prepare for errors.
\item BLP = confidentiality; Biba = integrity; Clark-Wilson = well-formed commercial transactions.
\item Hardware roots such as TPM/HSM protect keys and attest/perform cryptographic operations outside general software.
\item Modern architectures move trust boundaries; they do not eliminate them.
\item AES is the modern symmetric workhorse; asymmetric crypto solves key/signature problems; PKI binds keys to identity.
\item Crypto security is lifecycle security: algorithm, mode, randomness, key generation/storage/revocation, implementation, and agility.
\item Physical availability depends on power, cooling, environmental design, access control, and fire/life safety.
\end{itemize}

\section{Active-Recall Questions}
\begin{enumerate}
\item List the five ISO/IEC 19249 architectural principles and five design principles.
\item Compare Bell-LaPadula, Biba, and Clark-Wilson.
\item What do TPM attestation, binding, and sealing accomplish?
\item How do containers differ from VMs as an isolation boundary?
\item What changes in shared responsibility from IaaS toward SaaS?
\item Why is AES normally preferred to RSA for bulk data encryption?
\item What problem does Diffie-Hellman-Merkle solve, and what does it not provide by itself?
\item What does a digital signature provide, and what property does it not provide by itself?
\item Why must PKI validate more than the presence of a certificate?
\item Distinguish known-plaintext from chosen-plaintext and chosen-ciphertext attacks.
\item Name five observable artifacts a side-channel attack might use.
\item What are the five fire classes summarized in this chapter?
\end{enumerate}

\subsection*{Answer Key}
\begin{enumerate}
\item Architecture: domain separation, layering, encapsulation, redundancy, virtualization. Design: least privilege, attack-surface minimization, centralized parameter validation, centralized general security services, error/exception preparation.
\item BLP protects confidentiality; Biba integrity; Clark-Wilson commercial integrity through controlled transformations and separation of duties.
\item Attestation reports/measures trusted state; binding associates cryptographic protection with the platform; sealing releases data only when required platform state is present.
\item VMs emulate separate machines/guest OSs through a hypervisor; containers share the host kernel and isolate processes/resources.
\item The provider manages progressively more of the underlying stack, while the customer still retains responsibilities for identities, data, configuration, and appropriate use.
\item Symmetric AES is far faster for bulk data; RSA is asymmetric and used for public-key functions such as key establishment and signatures/encryption schemes.
\item It establishes a shared secret; unauthenticated exchange alone does not prove peer identity or prevent MITM.
\item Integrity, origin authentication, and support for nonrepudiation; not confidentiality by itself.
\item Validate issuer/chain, subject/name, time validity, usage constraints, and revocation/status, plus protect private keys.
\item Known plaintext is observed corresponding input/output; chosen plaintext lets the attacker request encryption of selected input; chosen ciphertext lets the attacker submit ciphertext to a decryption oracle.
\item Timing, cache behavior, power, electromagnetic emanations, error information.
\item A ordinary solids; B flammable liquids/gases; C energized electrical; D combustible metals; F/K cooking oils/greases.
\end{enumerate}

\section{Cornell Summary}
\begin{summarybox}{Domain 3 Summary}
Security architecture turns policy and risk into boundaries, mediated interfaces, trusted components, resilient services, and verifiable behavior. Formal models clarify the desired property; secure-design principles limit privilege and exposure; modern client/cloud/IoT/container/serverless/edge architectures require explicit trust boundaries; cryptography supplies confidentiality, integrity, identity, and nonrepudiation only when algorithms and keys are managed correctly; and physical facilities must preserve availability and evidence through access, power, cooling, environmental, and fire controls. The recurring CISSP question is whether the complete design -- including failure modes and operators -- continues to enforce its security requirements.
\end{summarybox}

\section{Final Domain Checklist}
\begin{itemize}[label=\(\square\)]
\item I can explain ISO/IEC 19249 and the chapter's related secure-design principles.
\item I can compare all major security models and their primary security property.
\item I can explain memory protection, TPMs, cryptoprocessors, and HSMs.
\item I can identify risks/controls for clients, servers, databases, ICS, cloud, distributed systems, IoT, microservices, containers, serverless, embedded, HPC, edge, and virtualization.
\item I can select symmetric/asymmetric/hash/PKI/signature mechanisms for the correct objective.
\item I can explain key and certificate lifecycle management.
\item I can recognize brute force, plaintext/ciphertext oracles, implementation, side-channel, timing, fault, MITM, credential, and ransomware threats.
\item I can design layered site/data-center/media/evidence/environmental/fire protection.
\end{itemize}
\vfill
{\footnotesize\color{Teal}Prepared as a study companion to Deane \& Kraus, \textit{The Official (ISC)\textsuperscript{2} CISSP CBK Reference}, 6th ed. Content is paraphrased and synthesized for study.}

\end{document}
